% Full written audits by descent_groups and dieudonne; root proof and source comparison.
\subsection{Scalar relations for residual hermitian groups}

Scalar presentations and binary chain arguments are classical.
For the commutative local case, Rogers--Schlichting
\cite[Theorems 1.2--1.3, pp.2--3]{C-RogersSchlichting2024}
give the chain lemma and the symmetric Grothendieck--Witt presentation.
For division rings with involution, Ranicki
\cite[Proposition 22.8, p.230]{C-Ranicki1992}
records a scalar presentation of the Witt group.
We supply the following complete proof with the rank retained,
both hermitian signs specified, and explicit operation extraction
in our equivariant setting. We claim no priority for these methods.

Let $D$ be a finite-dimensional division algebra over its center,
let $\sigma$ be an involution, and let $\epsilon\in\{1,-1\}$.
We assume characteristic different from two and write
\[
 S_\epsilon=\{a\in D^\times:\sigma(a)=\epsilon a\}.
\]
The center has degree at most two over its involution-fixed field;
hence $D$ is finite-dimensional over that fixed field as well.
The involution is an anti-automorphism, and forms are on right
$D$-spaces, with $h(v,w)=\epsilon\sigma(h(w,v))$.
All products in the following formulas retain their indicated order.

\begin{lemma}[Proved: a binary chain argument over an infinite fixed center]
\label{C:division_binary_chain}
Suppose the fixed field $F=Z(D)^\sigma$ is infinite.
Any two orthogonal bases of a nonsingular $\epsilon$-hermitian
space which admits such a basis are connected by a finite sequence
of changes affecting at most two basis vectors at a time.
\end{lemma}
\begin{proof}
We induct on the dimension $n$; dimension one permits scaling a
single vector. Let the two bases be $b_1,\ldots,b_n$ and
$c_1,\ldots,c_n$, with nonzero diagonal entries $a_i$ and $d_i$.
For a vector $w=\sum b_i x_i=\sum c_i y_i$, require that every
partial norm
\[
 \sum_{i=1}^j\sigma(x_i)a_ix_i,
 \qquad \sum_{i=1}^j\sigma(y_i)d_iy_i
 \quad(1\leq j\leq n)
\]
be nonzero. We justify simultaneous existence over $F$, not over
an unspecified extension. Choose an $F$-basis of $D$.
Each displayed expression is a polynomial map on the finite-dimensional
$F$-space underlying $D^n$. Each is nonzero as a polynomial map:
for the first family evaluate at $w=b_1$, and for the second at
$w=c_1$. For each expression choose one scalar coordinate which
is a nonzero polynomial. The product of these finitely many
polynomials is nonzero, and an infinite field has a rational
point where it does not vanish. This gives the required $w$.

Starting from the first basis, scale $b_1$ to $b_1x_1$.
This is legitimate since its norm is nonzero and hence $x_1\ne0$.
Successively replace the current first vector
$\sum_{i<j}b_ix_i$ and $b_j$ by an orthogonal basis of their plane
whose first vector is $\sum_{i\leq j}b_ix_i$.
Its norm is nonzero by construction. In a nonsingular two-dimensional
plane, a nondegenerate line has a nonsingular one-dimensional
orthogonal complement, so this replacement is possible.
The untouched vectors remain orthogonal to the plane.
After $n-1$ such changes the first vector is $w$.
Apply the same procedure to the second basis. The resulting bases
have common first vector $w$, and their other vectors are orthogonal
bases of the same nonsingular space $w^\perp$.
Induction connects those complements while fixing $w$.
Reversing the second sequence finishes the proof.
\end{proof}

\begin{theorem}[Proved: an explicit scalar-generator presentation]
\label{C:hermitian_scalar_presentation}
Unless $(D,\sigma,\epsilon)=(F,\mathrm{id},-1)$ for a commutative
field $F$, the split hermitian Grothendieck group
$GW^\epsilon(D,\sigma)$ has generators $[a]$, $a\in S_\epsilon$,
and precisely the following scalar relations:
\[
 [\sigma(c)ac]=[a]\qquad(c\in D^\times),                \tag{C.3.HS1}
\]
and, for $a,b\in S_\epsilon$ and $x\in D$ such that
$a'=a+\sigma(x)bx\ne0$, put
\[
 y=-a^{-1}\sigma(x)b,\qquad b'=b+\sigma(y)ay.
\]
Then $a',b'\in S_\epsilon$ and the second relation is
\[
                    [a]+[b]=[a']+[b'].                 \tag{C.3.HS2}
\]
These are explicit relations in at most three division-algebra
scalars; they contain no unknown higher-rank isometry predicate.
In the excluded alternating case the group is $\mathbb Z$ on
the standard alternating plane, with coordinate half the dimension.
The statement is a presentation, not a claim that its word problem
is decidable from arbitrary field arithmetic.
\end{theorem}
\begin{proof}
First every nonsingular form in the stated nonexceptional case
diagonalizes. Indeed, if $h(v,v)=0$ for every $v$, then polarization
gives $h(v,w)+h(w,v)=0$. As $h$ is nonsingular, for some fixed
nonzero $v$ its values $h(v,w)$ run through $D$.
It follows that $a+\epsilon\sigma(a)=0$ for all $a\in D$.
Taking $a=1$ gives $\epsilon=-1$, and then $\sigma=\mathrm{id}$.
An anti-automorphism equal to the identity forces $D$ commutative.
This is exactly the excluded case. Otherwise a vector of nonzero
norm splits off a nondegenerate line; induction on its orthogonal
complement diagonalizes. Thus the indicated generators suffice.

Relation (C.3.HS1) is scaling a basis vector.
For (C.3.HS2), start with Gram matrix $\operatorname{diag}(a,b)$
and take the new vectors $v_1=e_1+e_2x$ and $v_2=e_1y+e_2$.
Their cross-pairing is $ay+\sigma(x)b=0$.
Their norms are $a'$ and $b'$.
They are independent: if $v_2=v_1z$, their orthogonality gives
$a'z=0$, hence $z=0$, contradicting the second coordinate of $v_2$.
Therefore they are a basis of the original nonsingular plane,
so $b'\ne0$. The required symmetry of both new norms follows
from the hermitian identity. Both scalar relations are true.

Every change between two orthogonal bases of a plane follows
from these relations and permutation of summands. If its first
new vector has zero first coordinate, use permutation and scaling.
Otherwise scale the first old vector so the first new vector
has the form $e_1+e_2x$; (C.3.HS2) gives that vector and one
nonzero vector on its orthogonal complement. The target second
vector is a nonzero multiple of that complement vector, so
(C.3.HS1) finishes the change.

If $F=Z(D)^\sigma$ is infinite, the preceding lemma now implies
that any two diagonalizations of the same form give the same
sum of symbols modulo (C.3.HS1)--(C.3.HS2).
It follows that assigning the symbol sum of any diagonalization
is a well-defined additive inverse to the natural map from the
presented group to $GW^\epsilon(D,\sigma)$.
This also proves completeness for stable isometries in the group
completion: the common added diagonal sum cancels in the abelian
presentation. No cancellation theorem is needed for this step.

If $F$ is finite, then $D$ is finite and hence a finite field.
If $\sigma$ is nontrivial, all members of $S_\epsilon$ differ
by multiplication by a norm from $D^\times$ onto $F^\times$.
The first relation therefore identifies every generator with
one fixed generator, and the dimension homomorphism proves
that the presented group and the actual group are both $\mathbb Z$.
This includes skew-hermitian forms for the nontrivial involution.

For $\sigma=\mathrm{id}$ and $\epsilon=1$, there are two square
classes, represented by $1$ and a nonsquare $\nu$.
The finite-field form calculation already used in
Theorem~\ref{C:Om-finite-field-algorithm} says that
$\langle1,1\rangle\simeq\langle\nu,\nu\rangle$.
The plane argument above expresses this isometry using our scalar
relations, so $2[1]=2[\nu]$.
Consequently the presented group is a quotient of
$\mathbb Z^2/\langle2[1]-2[\nu]\rangle$.
Rank and determinant define its surjection onto the actual
$\mathbb Z\oplus\mathbb Z/2$, and identify that displayed upper
bound with the actual group. Both maps must be isomorphisms.
The commutative alternating case is the ordinary symplectic-basis
argument and gives one generator for each alternating plane.
This completes every case.
\end{proof}

\begin{corollary}[Proved reduction: scalar coordinates and all operations]
\label{C:WG_hermitian_scalar_operations}
In Proposition~\ref{C:field_WG_structure} and
Theorem~\ref{C:artinian_perfect_WG_reduction}, replace each residual
hermitian group by the scalar presentation of
Theorem~\ref{C:hermitian_scalar_presentation}, or by the alternating
plane coordinate in its exceptional case. Together with the integer
coordinates on non-self-dual pairs, this gives an explicit additive
presentation of the indicated split $WG$ group.

Suppose actual representative matrices, their endomorphism radicals,
and Krull--Schmidt decompositions of the input and resulting
tensor or exterior representations have been supplied by the preceding
classification constructions. Then every product and exterior
operation is expressed in the stated scalar generators by finitely
many linear algebra operations and tests of equality in the supplied
division algebras. No residual higher-rank hermitian-isometry oracle
is used. The result is a finite expression in the presentation;
deciding equality of arbitrary such expressions over arbitrary
coefficient fields is not asserted.
\end{corollary}
\begin{proof}
We give the extraction and lifting formulas, so the operations are
specified in these generators rather than left as an unspecified
form-class calculation. Use the notation $U,q_U,E_U,D_U,\sigma_U$
of Proposition~\ref{C:field_WG_structure}; the same notation applies
to an $R$-projective indecomposable in the Artinian theorem, with
all duals and forms taken over $R$.
Let $(M,q)$ be an actual symmetric representation and choose a
Krull--Schmidt decomposition. For a self-dual $U$ occurring with
multiplicity $r$, let $i_a:U\to M$, $1\leq a\leq r$, be the
specified inclusions. Its residual multiplicity Gram matrix is
\[
 H_U=\left(\overline{q_U^{-1}i_a^*q i_b}\right)_{a,b=1}^r
                  \in M_r(D_U).                         \tag{C.3.HS3}
\]
The bar means reduction of an actual endomorphism of $U$ modulo
$\operatorname{rad}(E_U)$. This is an
$\varepsilon_U$-hermitian matrix for $\sigma_U$.
The radical quotient of the additive category has no morphisms
between distinct indecomposable classes. Consequently the global
isomorphism $q:M\to M^*$ becomes an invertible block matrix there,
and every self-dual block (C.3.HS3) is nonsingular. Non-self-dual
blocks are exchanged by duality and contribute the integer $r$
on the hyperbolic generator for their unordered pair. This also
explains why nonorthogonal Krull--Schmidt inclusions are allowed:
the off-orbit components disappear in the category's radical
quotient, and the checked reduction and transfer theorem recovers
the form class from exactly these blocks.

In the nonexceptional case diagonalize (C.3.HS3) by choosing a
vector $v$ of nonzero norm, splitting its line, and repeating on
its orthogonal complement. Every step is finite linear algebra
over the division algebra, with the displayed multiplication order.
There is a finite search for such a vector from the supplied
finite-dimensional coefficient algebra. Indeed choose a basis of
$D_U$ over the coefficient field $k$. The $k$-coordinates of
$h(v,v)$ are polynomials, of degree at most two in each scalar
coordinate of $v$. At least one is a nonzero polynomial by the
diagonalization proof above. The grid
$\{0,1,-1\}^{r\dim_kD_U}$ therefore contains a vector of nonzero
norm, since $2\ne0$ and a polynomial of degree at most two in each
variable cannot vanish on that entire grid unless it is zero.
Test these finitely many vectors using the given multiplication
and involution tables. Once a vector is found, a linear basis change
puts it first, and subtracting its orthogonal projection from the
other vectors splits off its line. The resulting diagonal entries
$a_1,\ldots,a_r$ contribute
$[a_1]+\cdots+[a_r]$.
In the commutative alternating case use the ordinary symplectic
basis construction and output $r/2$. It is an integer because a
nonsingular alternating matrix has even size.

Conversely a scalar symbol $[a]$ has an explicit actual lift.
Choose $a_0\in E_U$ lifting $a$, and put
\[
 \widetilde a=(a_0+\varepsilon_U a_0^\dagger)/2,
 \qquad q_a=q_U\widetilde a,                             \tag{C.3.HS4}
\]
where $\dagger$ is the $q_U$-adjoint on $E_U$.
Then $\widetilde a^\dagger=\varepsilon_U\widetilde a$,
and its residue is the unit $a$, so it is a unit in $E_U$.
Thus $q_a$ is an actual nonsingular symmetric invariant form on
$U$. The alternating-plane generator lifts to $H(U)$, and the
integer generator for a non-self-dual pair also lifts to $H(U)$.
The checked nilpotent-radical lifting theorem makes each lift
independent of the choice of $a_0$ up to isometry. On a sum of
symbols take orthogonal sums of these lifts.

For products take the actual tensor representation and the tensor
Gram matrix. For $\lambda^j$ take the actual exterior representation
with its determinant Gram matrix; its entries are the corresponding
$j$-minors. Over a finite-dimensional coefficient algebra, a finite
projective module can be supplied as the image of an idempotent
$e$ on a finite free module. Its exterior power is the image of
$\bigwedge^j e$, whose entries are the same minors, so this rule
does not assume that the projective module is free or that $j!$
is invertible. Compute these objects before taking any radical
quotient. Use the supplied decomposition and (C.3.HS3) to extract
the result, followed by the finite diagonalization just described.
All objects and coefficient matrices in this operation are finite.
The radical need not be a tensor ideal, and no step treats it as one.

Different decompositions or lifts yield isometric actual input
forms by the proved transfer theorem. Tensor and exterior operations
preserve those isometries, and Theorem~\ref{C:hermitian_scalar_presentation}
identifies all resulting diagonalizations by (C.3.HS1)--(C.3.HS2).
This proves well-definedness in the stated presentation.
Products extend bilinearly to virtual classes; exterior coefficients
extend by $\lambda_t(x-y)=\lambda_t(x)\lambda_t(y)^{-1}$.
A specified coefficient uses only finitely many coefficients of
these series, and is computed by the same finite ring operations.
\end{proof}

\paragraph{Exact remaining scope.}
The residual hermitian groups and operation outputs now have the
explicit scalar relations (C.3.HS1)--(C.3.HS2), without an unknown
higher-rank isometry relation. The separate division-index and
primitive-idempotent calculation in the general perfect-field
construction remains arithmetic input. No uniform word-problem
algorithm over arbitrary fields, and no classification of every
remaining nonfield coefficient base, is concluded from this presentation.
